% Penggerak mandiri kumulatif untuk Logika Orde Pertama dan Dasar-Dasar
% Teori Model dalam pelokalan Bahasa Indonesia.

\documentclass[openany]{memoir}

% `cleveref` Open Logic belum memiliki opsi Indonesian. Babel tetap memuat
% Indonesian; `\ollanguage` bernilai english hanya selama cleveref dimuat.
\PassOptionsToPackage{english,indonesian}{babel}

\newcommand*{\olpath}{../..}
\newcommand*{\ollangid}{id}
\newcommand*{\ollanguage}{english}

\usepackage{mathpazo}
\usepackage{helvet}
\usepackage{courier}
\linespread{1.05}

\input{\olpath/sty/open-logic.sty}
\renewcommand*{\ollanguage}{indonesian}
\input{\olpath/sty/open-logic-defer.sty}

\includeenv{editorial}
\problemsperchapter
\let\cleardoublepage\clearpage
\pagestyle{plain}

% Dua rujukan bacaan lanjutan menunjuk ke bab 74 pembaca Inggris beku.
% Label sintetis mempertahankan nomor kanonik tanpa mengimpor bab tersebut.
\makeatletter
\hypertarget{ol-id-forward-inductive-definitions}{}
\protected@write\@auxout{}{\string\newlabel{mth:ind:idf:sec}{{74.4}{0}{Definisi Induktif}{ol-id-forward-inductive-definitions}{}}}
\protected@write\@auxout{}{\string\newlabel{mth:ind:idf:sec@cref}{{[section][4][74]74.4}{[1][0][]0}{}{}{}}}
\protected@write\@auxout{}{\string\newlabel{mth:ind:sti:sec}{{74.5}{0}{Induksi Struktural}{ol-id-forward-inductive-definitions}{}}}
\protected@write\@auxout{}{\string\newlabel{mth:ind:sti:sec@cref}{{[section][5][74]74.5}{[1][0][]0}{}{}{}}}
\makeatother

\begin{document}

\selectlanguage{indonesian}

% Pertahankan nomor bab pembaca lengkap: lima bab Logika Orde Pertama adalah
% bab 15--19; Dasar-Dasar Teori Model menjadi bab 20.
\setcounter{chapter}{14}

\olimport*[first-order-logic/introduction]{introduction}
\olimport*[first-order-logic/syntax-and-semantics]{syntax}
\olimport*[first-order-logic/syntax-and-semantics]{semantics}
\olimport*[first-order-logic/models-theories]{models-theories}
\olimport*[first-order-logic/beyond]{beyond}
\olimport*[model-theory/basics]{basics}

\bibliographystyle{\olpath/bib/natbib-oup}
\bibliography{\olpath/bib/open-logic}

\end{document}
